\item[\textbf{Exercise 2.1-5}]
Consider the problem of adding two $n$-bit binary integers $a$ and $b$, stored
in two $n$-element arrays $A[0 : n - 1]$ and $B[0 : n - 1]$, where each element
is either 0 or 1, $a = \sum_{i=0}^{n-1} A[i] \cdot 2^i$, and $b =
\sum_{i=0}^{n-1} B[i] \cdot 2^i$. The sum $c = a + b$ of the two integers should
be stored in binary form in an $(n + 1)$-element array $C[0 : n]$, where $c =
\sum_{i=0}^{n} C[i] \cdot 2^i$. Write a procedure \textsc{Add-Binary-Integers}
that takes as input arrays $A$ and $B$, along with the length $n$, and returns
array $C$ holding the sum.

\Solution

To determine the algorithm's behavior, consider the addition of two binary
sequences, $A$ and $B$, of length $n$:

\begin{enumerate}[
    style      = multiline,
    itemsep    = 2em,
]
    % Item 1
    \item $
    \begin{aligned}[t]
        A &= \bvec{0, 1, 1, 0, 1}, & n &= 5 \\
        B &= \bvec{1, 0, 0, 1, 0}, & n &= 5 \\
        \hline
        \text{SumBinaryLists}(A, B, n) \implies C &= \bvec{1, 1, 1, 1, 1, 0}, & n &= 6
    \end{aligned}
    $

    % Item 2
    \item $
    \begin{aligned}[t]
        A &= \bvec{1, 0, 0, 0, 1}, & n &= 5 \\
        B &= \bvec{1, 1, 0, 1, 0}, & n &= 5 \\
        \hline
        \text{SumBinaryLists}(A, B, n) \implies C &= \bvec{0, 0, 1, 1, 1, 0}, & n &= 6
    \end{aligned}
    $

    % Item 3
    \item $
    \begin{aligned}[t]
        A &= \bvec{0, 0, 0, 0, 0}, & n &= 5 \\
        B &= \bvec{0, 0, 0, 0, 0}, & n &= 5 \\
        \hline
        \text{SumBinaryLists}(A, B, n) \implies C &= \bvec{0, 0, 0, 0, 0, 0}, & n &= 6
    \end{aligned}
    $

    % Item 4
    \item $
    \begin{aligned}[t]
        A &= \bvec{1, 1, 1, 1, 1}, & n &= 5 \\
        B &= \bvec{1, 1, 1, 1, 1}, & n &= 5 \\
        \hline
        \text{SumBinaryLists}(A, B, n) \implies C &= \bvec{0, 1, 1, 1, 1, 1}, & n &= 6
    \end{aligned}
    $
\end{enumerate}

Based on the examples above and the exercise specifications, we can conclude
that a carry mechanism is required. The array $C$ requires an extra element (the
$(n)$-th position) to store the final overflowing value.

Here is the algorithm:

\begin{codebox}
\Procname{$\proc{SumBinaryLists}(A,B,n)$}
\li $C[0 \subarr n] \gets \{0\}$
\li $\id{carry} \gets 0$ 
\li \For $i \gets 1$ \To $n - 1$ \Do
\li     $\id{sum} \gets A[i] + B[i] + \id{carry}$
\li     $C[i] = \id{sum} \bmod 2$
\li     $\id{carry} \gets \lfloor \id{sum} / 2 \rfloor$
    \End
\li $C[n] \gets \id{carry}$
\li \Return $C$
\End
\end{codebox}
